\section{Polynomial sections on an elliptic curve}
\label{sec:elliptic}

Nine zero normal displacements impose zeros on the coefficients of a
normal polynomial. Dividing out these zeros puts its first $m+1$
coefficients in bundles of the same positive degree. We first make this
change intrinsically on any smooth cubic. We then choose cubics and
marked points for which theta functions give explicit bases of the
resulting spaces, ready for the collision and limit in
Section~\ref{sec:limits}.

In this section and Section~\ref{sec:limits}, $H^0$ denotes holomorphic
sections. The algebraic sections supplied by
Proposition~\ref{prop:normal-polynomial} belong to these spaces, so
excluding holomorphic solutions will also exclude the necessary normal
polynomials.

\subsection{Removing the nine prescribed coefficient zeros}

For this construction, let $d,m$ be any positive integers with $d>3m$.
Let $B$ be a smooth plane cubic, let $L=\mathcal O_B(1)$ and
$M=L^3$, and choose nine distinct points with divisor
$P=a_1+\cdots+a_9$. Write $\sigma_P$ for the canonical section of
$\mathcal O_B(P)$ with zero divisor $P$, and put
\[
 A=M\otimes\mathcal O_B(-P),\qquad
 T=L^{d-3m},\qquad u_0=\lfloor d/3\rfloor.
\]
Thus $\deg A=0$ and $\deg T=3(d-3m)>0$. Multiplication by
$\sigma_P$ defines a bundle map $A\to M$, written
\[
 w=\sigma_P v,
\]
which is invertible over $B\setminus P$. Define coefficient spaces
inside $H^0(B,T\otimes A^{m-j})$ by
\begin{equation}
 V_j=
 \begin{cases}
 H^0(B,T\otimes A^{m-j}),&0\leq j\leq m,\\[1mm]
 \sigma_P^{j-m}H^0(B,L^{d-3j}),&m<j\leq u_0.
 \end{cases}
\label{eq:coefficient-spaces}
\end{equation}
Negative powers denote dual powers. The target in the second line is
correct because
\[
 L^{d-3j}\otimes\mathcal O_B((j-m)P)=T\otimes A^{m-j}.
\]
Multiplication by $\sigma_P^{j-m}$ is injective, as can be checked on
$B\setminus P$. Let $\W$ be the space of polynomial sections on
$\Tot(A)$
\[
 \widetilde F(v)=\sum_{j=0}^{u_0}g_jv^j,\qquad g_j\in V_j,
\]
valued in the pullback of $T\otimes A^m=L^d(-mP)$.
It is a finite-dimensional space: its coefficient blocks are spaces of
sections on the compact curve $B$, and distinct powers of the fiber
coordinate make their sum direct.

The nine zero displacements are now built into $\W$. The remaining
points impose ordinary multiplicities at scalar fiber coordinate $1$,
as follows.

\begin{proposition}
\label{prop:necessary-W}
Suppose $r\geq10$, $d>3m$, and $\mathsf U_r(d,m)$ holds. Choose
$B,P,A$ and $\W$ as above, and fix a holomorphic frame of $A$ on
a coordinate disk in $B\setminus P$. For every choice of $q=r-9$
distinct points $b_1,\ldots,b_q$ in this disk, there is a nonzero
$\widetilde F\in\W$ of multiplicity at least $m$ at each point
$(b_\alpha,1)$ of $\Tot(A)$, where $1$ is the scalar coordinate in
the chosen frame.
\end{proposition}

\begin{proof}
Apply Proposition~\ref{prop:normal-polynomial} at the distinct points
$a_1,\ldots,a_9,b_1,\ldots,b_q$. Take zero normal displacements
at the $a_i$. At $b_\alpha$, take the image under multiplication by
$\sigma_P$ of the vector with scalar coordinate $1$ in the chosen
frame of $A$. The normal-polynomial proposition allows these arbitrary
normal vectors, so it supplies a nonzero
\[
 F(w)=\sum_{j=0}^{u_0}f_jw^j,\qquad
 f_j\in H^0(B,L^{d-3j}),
\]
with the required multiplicities.

By \eqref{eq:nine-zeros}, $f_j$ vanishes along $(m-j)P$ for $j<m$.
Consequently the expression
\[
 \widetilde F(v)=\sigma_P^{-m}F(\sigma_P v)
       =\sum_{j=0}^{u_0}f_j\sigma_P^{j-m}v^j
\]
extends across $P$. The prescribed zeros cancel its negative powers of
$\sigma_P$, and its coefficients have exactly the forms in
\eqref{eq:coefficient-spaces}. It therefore belongs to $\W$ and is
nonzero, since the fiber substitution is invertible on $B\setminus P$.

Near each $b_\alpha$, $w=\sigma_P v$ is a nonsingular change of
surface coordinates. The value-bundle map given by multiplication by
$\sigma_P^{-m}$ is also invertible there; in local frames it is
multiplication by a holomorphic unit. These two operations preserve
ordinary order of vanishing and take the displaced point to
$(b_\alpha,1)$.
\end{proof}

We next realize these coefficient spaces on explicit elliptic curves.
Their bases will retain the bundle multipliers, not merely their degrees;
those multipliers determine the Laurent exponents in the limit.

\subsection{Theta functions and plane cubics}

For $0<\tau<1$, let
\[
 X_\tau=\C^*/\tau^{\Z}
   \simeq\C/(2\pi i\Z+(\log\tau)\Z).
\]
This is a connected compact complex torus of dimension one. For
$n\in\Z$ and $\gamma\in\C^*$, define $D(n,\gamma)$ as the
holomorphic line bundle obtained from $\C^*\times\C$ by the action
generated by
\[
 (z,t)\longmapsto(\tau z,\gamma z^{-n}t).
\]
In its covering-space frame, sections are holomorphic functions with
\[
 f(\tau z)=\gamma z^{-n}f(z),
\]
and ordinary multiplication gives the tensor identifications
\[
 D(n,\gamma)\otimes D(n',\gamma')
   =D(n+n',\gamma\gamma').
\]

Define
\begin{equation}
 \Theta_\tau(z)
 =\prod_{p=0}^{\infty}(1-\tau^p z)
  \prod_{p=1}^{\infty}(1-\tau^p/z).
\label{eq:theta-product}
\end{equation}
We write $\Theta$ when $\tau$ is fixed. On every compact annulus,
the sums of the absolute deviations of the factors from $1$ are
bounded by geometric series. The products therefore converge locally
uniformly on $\C^*$, and reindexing gives
\[
 \Theta(\tau z)=-z^{-1}\Theta(z).
\]
Its zeros are precisely the points of $\tau^{\Z}$, all simple:
exactly one factor vanishes at each, while the remaining product is
nonzero. Thus $\Theta(z/a)$ is a section of $D(1,-a)$ with divisor
the single point represented by $a\in\C^*$. Products of $n$ such
sections, with the parameters having product $(-1)^n\gamma$, show
that $D(n,\gamma)$ has degree $n$ for $n>0$. Tensor products and
duals give the formula for every integer $n$; in degree zero, use
$D(0,\gamma)=D(1,\gamma)\otimes D(1,1)^{-1}$.

The multiplicative theta convention and residue-class bases are classical;
see Rosengren~\cite[arXiv version, Section~1.2 and Exercise~1.6.5]{Rosengren2020}.
We include the proof to fix all multipliers and normalizations.

\begin{lemma}
\label{lem:torus-sections}
For $n>0$ and $\gamma\in\C^*$, the space
$H^0(X_\tau,D(n,\gamma))$ has dimension $n$. The Laurent
coefficients of a section $f(z)=\sum_{K\in\Z}c_Kz^K$ satisfy
\begin{equation}
\begin{split}
 c_{K+n}&=\gamma^{-1}\tau^Kc_K,\\
 c_{K+np}&=c_K\gamma^{-p}
       \tau^{pK+np(p-1)/2}\qquad(p\in\Z).
\end{split}
\label{eq:laurent-recurrence}
\end{equation}
One coefficient in each residue class modulo $n$ may be prescribed
arbitrarily. Taking one nonzero starting coefficient at a time gives
a basis supported in the individual residue classes.
\end{lemma}

\begin{proof}
A holomorphic function on $\C^*$ has a Laurent expansion convergent
on compact annuli. Comparing coefficients in its functional equation
gives $\tau^Kc_K=\gamma c_{K+n}$; iteration in either direction
proves \eqref{eq:laurent-recurrence}.

Conversely, prescribe one coefficient per residue class and extend
by the recurrence. For a fixed starting exponent $K$, the terms on
any compact annulus are bounded by
\[
 C\exp(-c p^2+C'|p|),\qquad c>0,
\]
because $n>0$ and $\log\tau<0$. The resulting Laurent series
converges normally and satisfies the functional equation. There are
exactly $n$ free coefficients, and uniqueness of Laurent expansions
proves the basis assertion.
\end{proof}

\begin{lemma}
\label{lem:cubic-embedding}
For every $\gamma\in\C^*$, the complete space of sections of
$D(3,\gamma)$ embeds $X_\tau$ into $\Pp^2$ as a smooth irreducible
plane cubic, with $\mathcal O_{\Pp^2}(1)$ pulling back to
$D(3,\gamma)$.
\end{lemma}

\begin{proof}
Let $Q$ be an effective divisor of degree $t\in\{1,2\}$. A product
of $t$ theta factors, using lifts of its points with their
multiplicities, is a section with divisor exactly $Q$. Division by
this product identifies sections of $D(3,\gamma)$ vanishing along
$Q$ with sections of some $D(3-t,\gamma')$. Their dimension is
$3-t$ by Lemma~\ref{lem:torus-sections}.

For $t=1$, this proves that the three-dimensional section space has
no base points. For two distinct points it shows that the associated
map separates them. For $Q=2p$, the dimension drop from sections
vanishing at $p$ to those vanishing twice supplies a section of order
exactly one at $p$. Its ratio with a section nonzero there has nonzero
derivative. The map is therefore an injective holomorphic immersion,
and compactness makes it an embedding.

By Chow's Theorem~\cite{Chow1949}, the closed analytic image is
algebraic. It is smooth and irreducible because it is an embedded
connected compact Riemann surface. Its hyperplane bundle pulls back
to $D(3,\gamma)$, of degree three, so the image is a plane cubic.
\end{proof}

\subsection{Choosing the cubic and marked points}

We now fix a nonsquare integer $r\geq10$ and a hypothetical universal
system $\mathsf U_r(d,m)$ with $d/m\leq\sqrt r$. The rationality
of $d/m$ and Lemma~\ref{lem:cubic-low-degree} give
$3<d/m<\sqrt r$. Set
\begin{equation}
\begin{gathered}
 q=r-9,\qquad
 \rho=3(d/m-3),\qquad \rho_*=3(\sqrt r-3),\\
 0<\rho<\rho_*,\qquad
 2\rho_*+\frac{\rho_*^2}{9}=q.
\end{gathered}
\label{eq:rho}
\end{equation}
Put
\[
 h=3(d-3m)=m\rho>0,\qquad u_0=\lfloor d/3\rfloor\geq m.
\]

Set $x_0=1/2$ and choose a rational
number $\Delta$ such that
\begin{equation}
 0<\Delta<x_0,\qquad
 \rho+\frac{\Delta\rho}{9}<\rho_*.
\label{eq:delta}
\end{equation}
The strict slope gap leaves an interval of such choices.
Rationality of $\Delta$ will keep the Laurent exponent
intervals away from integral endpoints.

Choose nine distinct real numbers $x_i$ with
\begin{equation}
\begin{gathered}
 0<x_i<x_0,\qquad \sum_{i=1}^{9}(x_0-x_i)=\Delta,\\
 \eta=\frac{\sum_{i=1}^{9}x_i+\rho_*}{3},\qquad
 z_i=\tau^{x_i}\ (1\leq i\leq9),\qquad z_0=\tau^{x_0}.
\end{gathered}
\label{eq:marked-parameters}
\end{equation}
For example, take $x_i=x_0-i\Delta/45$. The quantities
$\rho,\Delta,x_i,\eta$ can be kept unchanged under a common
multiple of $(d,m)$, whereas $h,u_0$ and the coefficient spaces
depend on the actual pair. We keep this pair fixed while taking the
collision on each curve and the subsequent limit in $\tau$.

For each $0<\tau<1$, use Lemma~\ref{lem:cubic-embedding} to identify
$X_\tau$ with a plane cubic $B$, with the holomorphic identification
\[
 L=\mathcal O_B(1)=D(3,\tau^\eta),\qquad M=L^3.
\]
Let $a_i$ be the point represented by $z_i$ and put
$P=a_1+\cdots+a_9$. These points and the point represented by $z_0$
are distinct: their exponents lie in $(0,1)$, and two lifts represent
the same point only if their exponents differ by an integer.

In the covering-space frame, the section
\[
 \Pi(z)=\prod_{i=1}^{9}\Theta(z/z_i)
\]
has divisor $P$ and gives an identification
\[
 \mathcal O_B(P)=D\left(9,-\tau^{\sum_i x_i}\right)
\]
under which the canonical section $\sigma_P$ is represented by
$\Pi(z)$. Thus the bundles already used in the construction are
\[
 A=M\otimes\mathcal O_B(-P),\qquad T=L^{d-3m}.
\]
Since $3\eta=\sum_i x_i+\rho_*$, we have
\[
 A=D(0,-\tau^{\rho_*}),\qquad
 \deg T=h,\qquad
 T\otimes A^m=L^d\otimes\mathcal O_B(-mP).
\]
These identifications use the covering-space frames and their tensor
products. We retain the displayed powers of $\tau$ as written,
without replacing the multiplier presentations by isomorphic ones.

A cubic equation identifies the normal bundle of $B$ with $M=L^3$.
Its differential at each point surjects onto the normal fiber, so
every scalar specified in the chosen covering-space frame is an
allowed normal displacement in
Proposition~\ref{prop:normal-polynomial}. Rescaling the cubic equation
does not restrict those choices. All constructions are made for each
fixed $\tau$; their algebraic dependence on $\tau$ is not required.

Denote the resulting polynomial space by $\W_\tau$. By
Lemma~\ref{lem:torus-sections}, its first $m+1$ coefficient blocks
have degree and dimension $h$. For $j>m$, the dimension is
$3(d-3j)$ when $d-3j>0$, and is one when $d-3j=0$. In the last
case the source is $H^0(B,\mathcal O_B)=\C$, since $B$ is compact
and connected.

Near the point represented by $(z,v)=(z_0,1)$, use coordinates
\begin{equation}
 z=z_0e^x,\qquad v=e^y.
\label{eq:local-chart}
\end{equation}
Choose the $x$-disk small enough to map injectively to $B$ and avoid
$P$, and the $y$-disk small enough that $e^y$ is injective. In the
fixed covering-space frame of $T\otimes A^m$, the scalar expression
is
\[
 \widetilde F(x,y)=\sum_{j=0}^{u_0}g_j(z_0e^x)e^{jy}.
\]
Here $g_j(z)$ denotes the holomorphic function in its covering-space
frame; all summands take values in the same bundle.

For every fixed $\tau\in(0,1)$ and every $q$ distinct positions
$\xi_1,\ldots,\xi_q$ sufficiently near zero, apply
Proposition~\ref{prop:necessary-W} in this disk and the fixed frame
of $A$. It gives a nonzero $\widetilde F\in\W_\tau$ of multiplicity
at least $m$ at each $(x,y)=(\xi_\alpha,0)$, since $v=1$ corresponds
to $y=0$. The existence is pointwise in $\tau$ and in the tuple of
positions; the collision will use no simultaneous choice of sections.
